\documentclass[12pt]{article}

% Idioma e codificação
\usepackage[brazil]{babel}
\usepackage[utf8]{inputenc}
\usepackage[T1]{fontenc}

% Matemática e listas
\usepackage{amsmath, amssymb}
\usepackage{enumitem}

% Layout e cores
\usepackage{geometry}
\usepackage{graphicx}
\usepackage{xcolor}

% Cores institucionais
\definecolor{maincolor}{RGB}{0,100,50}
\definecolor{forestgreen}{RGB}{0,100,50}

% Cabeçalho e rodapé
\usepackage{fancyhdr}

% Blocos
\usepackage[most]{tcolorbox}
\newtcolorbox{block}[1]{%
  colback=green!5,
  colframe=maincolor,
  coltitle=white,
  title=#1,
  fonttitle=\bfseries,
  boxrule=1.5pt,
  arc=3mm,
  left=6pt,
  right=6pt,
  top=6pt,
  bottom=6pt
}

% Margens
\geometry{margin=2.5cm}

% Fancyhdr
\pagestyle{fancy}
\fancyhf{}

\renewcommand{\headrulewidth}{2pt}
\renewcommand{\footrulewidth}{2pt}

\renewcommand{\headrule}{\hbox to\headwidth{%
  \color{forestgreen}\leaders\hrule height \headrulewidth\hfill}}

\renewcommand{\footrule}{\hbox to\headwidth{%
  \color{forestgreen}\leaders\hrule height \footrulewidth\hfill}}

% Altura do cabeçalho (corrigida)
\setlength{\headheight}{52pt}
\addtolength{\topmargin}{-20pt}

% Cabeçalho
\fancyhead[L]{%
  \includegraphics[height=1.2cm]{Universidade_de_Rio_Verde.jpg}
}

\fancyhead[R]{%
  \raggedleft
  \textcolor{forestgreen}{%
    \small
    \textbf{Lista de Exercícios}\\
    Me.~Sergio Souza Novak\\
    \textit{Engenharia da Qualidade e Confiabilidade}
    }
    }
    
    % Rodapé
    \fancyfoot[C]{\textcolor{forestgreen}{\thepage}}
    
    % =========================
    \begin{document}
    % =========================
    
    \vspace*{0.2cm}
    \begin{center}
      {\Large \textbf{Atividade Avaliativa 1}}\\[0.2cm]
      {\large Engenharia da Qualidade e Confiabilidade}\\[0.2cm]
      Tema: Confiabilidade de Serviços, DevOps e SRE
    \end{center}
    
    \vspace{0.5cm}
    
    % \noindent
    % \textbf{Aluno(a):} \rule{10cm}{0.4pt} \hfill
    % \textbf{Data:} \rule{3cm}{0.4pt}
    
    % \vspace{0.5cm}
    
    
    \textbf{1.} Explique o que é \textbf{Engenharia da Qualidade} e qual é seu principal objetivo dentro das organizações modernas.\\[0.3cm]
    \rule{\textwidth}{0.4pt}\\
    \rule{\textwidth}{0.4pt}\\
    \rule{\textwidth}{0.4pt}
    
    \vspace{0.3cm}
    
    
    
    \textbf{2.} Descreva o conflito estrutural entre as equipes de \textbf{Desenvolvimento (Dev)} e \textbf{Operações (Ops)} no modelo tradicional. Explique por que esse conflito motivou o surgimento de abordagens como DevOps e SRE.\\[0.3cm]
    \rule{\textwidth}{0.4pt}\\
    \rule{\textwidth}{0.4pt}\\
    \rule{\textwidth}{0.4pt}\\
    \rule{\textwidth}{0.4pt}
    
\vspace{2cm}
    
    
    Assinale apenas \textbf{uma alternativa correta} em cada questão.
    
    \begin{enumerate}[label=\textbf{\arabic*.}]
    \item O principal foco da equipe de \textbf{Operações (Ops)} é:
    \begin{itemize}
      \item ( \ ) Lançar novas funcionalidades rapidamente
    \item ( \ ) Aumentar a complexidade do sistema
    \item ( \ ) Manter a estabilidade e a confiabilidade do serviço
    \item ( \ ) Desenvolver novas features para o produto
\end{itemize}

\item Segundo a definição do Google, \textbf{Site Reliability Engineering (SRE)} é:
\begin{itemize}
    \item ( \ ) Uma ferramenta de automação de deploy
    \item ( \ ) Um modelo de terceirização de operações
    \item ( \ ) O resultado de aplicar engenharia de software à operação de sistemas
    \item ( \ ) Um sinônimo direto de DevOps
\end{itemize}

\item No modelo SRE do Google, nenhuma equipe deve gastar mais de 50\% do tempo em:
\begin{itemize}
    \item ( \ ) Reuniões de planejamento
    \item ( \ ) Desenvolvimento de novas funcionalidades
    \item ( \ ) Atividades operacionais manuais
    \item ( \ ) Monitoramento automático
\end{itemize}

\item O uso excessivo de \textbf{feature flags} pode resultar em:
\begin{itemize}
    \item ( \ ) Redução completa da complexidade do sistema
    \item ( \ ) Eliminação de falhas em produção
    \item ( \ ) Aumento da complexidade lógica e de configuração
    \item ( \ ) Fim da necessidade de testes
\end{itemize}
\end{enumerate}

\vspace{0.3cm}


Relacione corretamente os itens da \textbf{Coluna A} com a \textbf{Coluna B}.

\vspace{0.3cm}

\subsection*{5. Engenharia da Qualidade, Operação e SRE}

\noindent
\begin{tabular}{p{6cm} p{8cm}}
\textbf{Coluna A} & \textbf{Coluna B} \\[0.2cm]
( 1 ) Engenharia da Qualidade & ( \ \ ) Prática que busca confiabilidade por meio de automação e métricas \\
( 2 ) Sysadmin tradicional & ( \ \ ) Foco em planejamento, controle e melhoria de processos \\
( 3 ) SRE & ( \ \ ) Operação baseada em tarefas manuais e intervenção humana \\
( 4 ) Postmortem & ( \ \ ) Análise de incidentes sem atribuição de culpa \\
\end{tabular}

\vspace{0.6cm}

\subsection*{6. DevOps e SRE}

\noindent
\begin{tabular}{p{6cm} p{8cm}}
\textbf{Coluna A} & \textbf{Coluna B} \\[0.2cm]
( 1 ) DevOps & ( \ \ ) Uso formal de SLI, SLO e Error Budget \\
( 2 ) SRE & ( \ \ ) Cultura de colaboração entre desenvolvimento e operações \\
( 3 ) DevOps & ( \ \ ) Abordagem mais filosófica e menos prescritiva \\
( 4 ) SRE & ( \ \ ) Disciplina de engenharia com limites claros de operação \\
\end{tabular}

\vspace{0.8cm}


% =========================
\end{document}
% =========================
